%% kapa: exam=1 round=34 year=2023 period=1 form=A subject=real_estate subject_ko=부동산학원론 seq=21 number=101
%% source: Q-Net 감정평가사 기출문제 (cst00309) artlSeq=5211789 file=raw/1008/5211789/2023년도 제34회 감정평가사 1차시험 1교시.pdf
%% conversion: kordoc 4.24.51 -> tools/gen1.py (auto); review=proofread-sonnet; flags=none
\begin{question}{101}
주택금융의 상환방식에 관한 설명으로 옳지 않은 것은?
\begin{choices}
  \item 만기일시상환방식은 대출만기 때까지는 원금상환이 전혀 이루어지지 않기에 매월 내는 이자가 만기 때까지 동일하다.
  \item 원금균등분할상환방식은 대출 초기에 대출원리금의 지급액이 가장 크기에 차입자의 원리금지급 부담도 대출 초기에 가장 크다.
  \item 원리금균등분할상환방식은 매기의 대출원리금이 동일하기에 대출 초기에는 대체로 원금상환 부분이 작고 이자지급 부분이 크다.
  \item 점증상환방식은 초기에 대출이자를 전부 내고, 나머지 대출원금을 상환하는 방식으로 부의 상환(negative amortization)이 일어날 수 있다.
  \item 원금균등분할상환방식이나 원리금균등분할상환방식에서 거치기간을 별도로 정할 수 있다.
\end{choices}
\end{question}
\answer{4}
